\chapter{Simulación Circuital en LTspice}
\label{chap:simulacion}

Con el fin de validar exhaustivamente el diseño analítico y contrastar el comportamiento dinámico de la etapa amplificadora,
se implementó el circuito en el entorno de simulación LTspice (archivo esquemático \texttt{ej\_Amp.SurtidorComun1.asc}).
Se efectuó el análisis del punto estático de polarización continua mediante la directiva \lstinline[style=ltspice]{.op} y
el análisis temporal en régimen permanente periódico mediante \lstinline[style=ltspice]{.tran}.

\section{Directivas de Simulación y Modelo SPICE del Transistor 2N7000}

Para modelar con exactitud el comportamiento físico del dispositivo ensayado en laboratorio, se integró el subcircuito
SPICE oficial del fabricante ON Semiconductor, cuyos parámetros de transconductancia y tensión umbral fueron configurados
en correspondencia estricta con los valores experimentales extraídos en el Capítulo~\ref{chap:caracterizacion_mosfet}:
$\mathrm{VTO} = 1.60104\V$ y $\mathrm{KP} = 0.053905\A/\V^2$ ($53.905\mA/\V^2$).

En el listado~\ref{list:spice_netlist_tp5} se transcribe la definición del subcircuito y las directivas de control
configuradas en el simulador.

\begin{lstlisting}[style=ltspice, caption={Subcircuito del transistor 2N7000 y directivas de simulación en LTspice.}, label={list:spice_netlist_tp5}]
* Modelo SPICE de alta precision para MOSFET Canal N 2N7000
.SUBCKT 2N7000 1 2 3
* Nodo 1 -> Drenador, Nodo 2 -> Compuerta, Nodo 3 -> Surtidor
M1 9 7 8 8 MM L=100u W=100u
.MODEL MM NMOS LEVEL=1 IS=1e-32 VTO=1.60104 LAMBDA=0 KP=0.053905
+ CGSO=1.79115e-07 CGDO=1.0724e-11
RS 8 3 1.10523
D1 3 1 MD
.MODEL MD D IS=2.71011e-10 RS=0.0140826 N=1.5 BV=60 IBV=1e-05
+ EG=1.16084 XTI=3.00131 TT=0 CJO=3.41211e-11 VJ=4.67429 M=0.899864
RDS 3 1 2.4e+11
RD 9 1 0.0001
RG 2 7 2.18034
D2 4 5 MD1
.MODEL MD1 D IS=1e-32 N=50 CJO=7.93181e-11 VJ=0.643298 M=0.9
D3 0 5 MD2
.MODEL MD2 D IS=1e-10 N=0.400165 RS=3.00002e-06
RL 5 10 1
FI2 7 9 VFI2 -1
VFI2 4 0 0
EV16 10 0 9 7 1
CAP 11 10 1.58786e-10
FI1 7 9 VFI1 -1
VFI1 11 6 0
RCAP 6 10 1
D4 0 6 MD3
.MODEL MD3 D IS=1e-10 N=0.400165 RS=0 BV=1000 IBV=1mA
.ENDS 2N7000

* Fuente senoidal de entrada: f = 1 kHz, amplitud 45 mVp (90 mVpp)
Vi in 0 SINE(0 45mV 1kHz)
.tran 50m
.op
\end{lstlisting}

\section{Simulación del Punto de Operación en Continua (.op)}

A partir de la ejecución del comando \lstinline[style=ltspice]{.op}, los datos numéricos extraídos del archivo binario
\texttt{ej\_Amp.SurtidorComun1.raw} quedan plasmados en el Cuadro~\ref{tab:resultados_op_ltspice}.

\begin{table}[!ht]
  \centering
  \resizebox{0.92\textwidth}{!}{
  \begin{tabular}{|l|l|c|c|}
    \hline
    \textbf{Variable SPICE} & \textbf{Parámetro Circuital} & \textbf{Valor Simulado} & \textbf{Unidad} \\
    \hline
    \lstinline[style=ltspice]{V(n001)} & Tensión continua de alimentación $V_{DD}$ & $15.0000$ & $\V$ \\
    \lstinline[style=ltspice]{V(g)} & Tensión continua de compuerta $V_G$ & $3.2500$ & $\V$ \\
    \lstinline[style=ltspice]{V(s)} & Tensión continua de surtidor $V_S$ & $1.4580$ & $\V$ \\
    \lstinline[style=ltspice]{V(d)} & Tensión continua de drenador $V_D$ & $8.3905$ & $\V$ \\
    \lstinline[style=ltspice]{V(g)-V(s)} & Tensión compuerta-surtidor de reposo $V_{GSQ}$ & $1.7920$ & $\V$ \\
    \lstinline[style=ltspice]{V(d)-V(s)} & Tensión drenador-surtidor de reposo $V_{DSQ}$ & $6.9325$ & $\V$ \\
    \lstinline[style=ltspice]{I(RD)} & Corriente continua de drenador $I_{DQ}$ & $0.9720$ & $\mA$ \\
    \lstinline[style=ltspice]{I(RS)} & Corriente continua de surtidor $I_{SQ}$ & $0.9720$ & $\mA$ \\
    \lstinline[style=ltspice]{I(R1)} & Corriente en la rama inferior del divisor & $25.00$ & $\uA$ \\
    \lstinline[style=ltspice]{I(R2)} & Corriente en la rama superior del divisor & $25.00$ & $\uA$ \\
    \lstinline[style=ltspice]{Ix(m1:2)} & Corriente continua de compuerta $I_G$ & $0.00$ & $\uA$ \\
    \lstinline[style=ltspice]{I(VDD)} & Consumo total de corriente desde la fuente & $-0.9970$ & $\mA$ \\
    \hline
  \end{tabular}
  }
  \caption{Resultados numéricos estáticos extraídos de la simulación .op en LTspice.}
  \label{tab:resultados_op_ltspice}
\end{table}

En el Cuadro~\ref{tab:comparacion_op_tp5} se confrontan los resultados teóricos con los datos simulados por LTspice.

\begin{table}[!ht]
  \centering
  \resizebox{0.88\textwidth}{!}{
  \begin{tabular}{|c|c|c|c|}
    \hline
    \textbf{Magnitud de Reposo} & \textbf{Teórico Normalizado} & \textbf{Simulación LTspice} & \textbf{Diferencia Relativa} \\
    \hline
     $V_G\ [\V]$ & $3.250$ & $3.250$ & $0.00\%$ \\
     $V_S\ [\V]$ & $1.459$ & $1.458$ & $-0.07\%$ \\
     $V_D\ [\V]$ & $8.386$ & $8.390$ & $+0.05\%$ \\
     $V_{GSQ}\ [\V]$ & $1.791$ & $1.792$ & $+0.06\%$ \\
     $V_{DSQ}\ [\V]$ & $6.927$ & $6.932$ & $+0.07\%$ \\
     $I_{DQ}\ [\mA]$ & $0.9727$ & $0.9720$ & $-0.07\%$ \\
     $I_G\ [\uA]$ & $0.000$ & $0.000$ & $0.00\%$ \\
    \hline
  \end{tabular}
  }
  \caption{Confrontación entre el cálculo analítico normalizado y la simulación .op en LTspice.}
  \label{tab:comparacion_op_tp5}
\end{table}

La discrepancia entre la formulación analítica y la simulación es prácticamente imperceptible (inferior al $0.07\%$),
lo que verifica la estricta consistencia entre las ecuaciones deducidas y los algoritmos numéricos de SPICE.

\section{Análisis de Formas de Onda Dinámicas y Ganancia (.tran)}

Para el análisis temporal se inyectó una excitación senoidal pura de $f = 1\kHz$ con amplitud de $45\mV_p$ ($90\mVpp$).
Se integró durante un lapso total de $50\ms$ para superar holgadamente el régimen transitorio impuesto por los
capacitores de acoplo y desacoplo ($C_1 = C_2 = C_3 = 100\uF$). En la figura~\ref{fig:sim_waveforms} se presentan las
formas de onda de tensión en régimen permanente periódico entre $48\ms$ y $50\ms$.

\begin{figure}[!ht]
  \centering
  \begin{tikzpicture}
  \pgfplotsset{set layers}
  \begin{axis}[
    scale only axis,
    width=0.68\textwidth,
    height=6.5cm,
    xlabel={Tiempo $[\mathrm{ms}]$},
    ylabel={Tensión de entrada $v_i\ [\mathrm{mV}]$},
    ylabel style={color=blue},
    yticklabel style={color=blue},
    ymin=-60, ymax=60,
    xmin=48.0, xmax=50.0,
    grid=both,
    grid style={line width=.1pt, draw=gray!30},
    major grid style={line width=.2pt, draw=gray!50},
    axis y line*=left,
    every axis plot/.append style={very thick}
  ]
    \addplot[color=blue, smooth] table [x=time, y=vi] {figures/waveforms_data.dat};
  \end{axis}
  \begin{axis}[
    scale only axis,
    width=0.68\textwidth,
    height=6.5cm,
    ylabel={Tensión de salida $v_L\ [\mathrm{mV}]$},
    ylabel style={color=red!80!black},
    yticklabel style={color=red!80!black},
    ymin=-2500, ymax=2500,
    xmin=48.0, xmax=50.0,
    axis y line*=right,
    axis x line=none,
    every axis plot/.append style={very thick}
  ]
    \addplot[color=red!80!black, dashed, smooth] table [x=time, y=vl] {figures/waveforms_data.dat};
  \end{axis}
\end{tikzpicture}

  \caption{Formas de onda senoidales simuladas en LTspice: tensión de entrada $v_i(t)$ y tensión de salida sobre la carga $v_L(t)$.}
  \label{fig:sim_waveforms}
\end{figure}

Del análisis de los datos temporales del régimen permanente periódico se constata:
\begin{itemize}
  \item Tensión de entrada pico a pico: $\hat{v}_{i,\mathrm{pp}} = 89.97\mVpp$ (amplitud pico de $44.98\mV_p$).
  \item Tensión de salida sobre la carga: $\hat{v}_{L,\mathrm{pp}} = 3865.12\mVpp \approx 3.865\Vpp$ (amplitud pico de $1.932\V_p$).
  \item Ganancia de tensión de la etapa:
    \begin{equation}
      A_v = \frac{v_L}{v_i} = \frac{-1932.56\mV_p}{44.98\mV_p} \approx -42.96
      \label{eq:Av_simulada}
    \end{equation}
  \item Desfase angular: $180^\circ$ (inversión de fase perfecta respecto a la compuerta).
\end{itemize}

La ganancia simulada ($|A_v| = 42.96$) coincide con el cálculo analítico de pequeña señal
($|A_v| = 42.91$), presentando un error despreciable de $+0.12\%$. Las formas de onda son puramente senoidales,
sin ningún tipo de asimetría o recorte dinámico.
